\documentclass[../main.tex]{subfiles}

\begin{document}

\chapter*{Résumé}
\addcontentsline{toc}{chapter}{Résumé}

Ce projet de fin d'études a été réalisé au sein du Pôle Business Développement
d'Atlantic Re, réassureur international filiale du Groupe CDG, dans le cadre d'une
mission initiale de modélisation stratégique de l'expansion africaine à horizon 2030.
L'objectif premier était de construire un modèle de priorisation des marchés africains
Vie et Non-Vie, fondé sur un pipeline data science multicritère.

Au cours de la phase d'immersion métier, nous avons détecté un problème structurel
qui n'entrait pas dans le périmètre initial : l'absence d'outil de consultation et
d'exploitation des données internes, rendant difficile pour les équipes du département
l'analyse des portefeuilles existants et, par conséquent, la fiabilité du modèle
externe que nous construisions. Face à ce constat, nous avons proposé d'élargir
notre mission à la digitalisation, en développant un écosystème digital servant de
portail de gouvernance de la donnée interne.

Le projet couvre ainsi deux axes complémentaires : d'une part, l'audit de la base
de données métier et le développement de l'écosystème digital ; d'autre part, la
construction du pipeline data science (collecte multi-sources, imputation ML
hiérarchique, scoring multicritère SCAR) aboutissant à un classement opérationnel
des marchés africains pour Atlantic Re.

\textbf{Mots-clés :} réassurance, Afrique, data science, machine learning, scoring
multicritère, gouvernance des données, digitalisation, expansion stratégique.

\clearpage

% ============================================================
%  ABSTRACT — ANGLAIS
% ============================================================

\end{document}
